\label{ch:sterile-observables}

\section{Named sterile presets}
\label{sec:sterile-presets-table}

\modulehead{config/presets.py}{\code{OSCILLATION\_PRESETS} entries carrying
a \code{theta14\_deg} key: named, literature-justified 3+1 parameter sets,
dispatched to \pyclass{PMNS\_sterile} automatically.}

\begin{implbox}
A preset is a \emph{sterile} preset purely by having a \code{theta14\_deg}
key; \pyfunc{PropagationConfig.oscillation\_parameters\_from\_preset}
dispatches on that key alone to build a \pyclass{PMNS\_sterile} instead of
a \pyclass{PMNS\_SM}, with every other field (the SM angles,
$\Delta m^2_{21}/\Delta m^2_{3\ell}$) inherited from the shared NuFIT 5.2
baseline~\parencite{Esteban2020NuFit,NuFit52}. As with the NSI presets of
the companion document, these are named, traceable benchmark points rather
than a single global best fit: current 3+1 data is in tension between
different anomalies (\cref{ch:sterile-theory}), so no universally preferred
point exists.

\begin{center}
\small
\begin{tabular}{@{}ll@{}}
\toprule
Preset & Parameters and origin \\
\midrule
\parbox[t]{3.4cm}{\code{sterile\_\allowbreak{}3p1\_\allowbreak{}null\_\allowbreak{}mixing}} &
  \parbox[t]{9.2cm}{$\theta_{14}=\theta_{24}=\theta_{34}=0$,
  $\Delta m^2_{41}=1.0\ \mathrm{eV^2}$. Null-hypothesis reference: with
  every sterile angle at zero, \pyclass{PMNS\_sterile} is numerically
  identical to the plain 3-flavour SM run through the 4-flavour machinery
  -- validated to $\sim10^{-16}$ in \cref{sec:results-sterile1}.} \\[26pt]
\parbox[t]{3.4cm}{\code{sterile\_\allowbreak{}3p1\_\allowbreak{}bestfit\_\allowbreak{}giunti2017}} &
  \parbox[t]{9.2cm}{$\theta_{14}=8.5^\circ$, $\theta_{24}=7.5^\circ$,
  $\theta_{34}=0$, $\Delta m^2_{41}=1.7\ \mathrm{eV^2}$
  ($\sin^22\theta_{14}=0.085$, $\sin^22\theta_{24}=0.068$). Global 3+1 fit
  combining the reactor and gallium anomalies with solar data
  \parencite{GariazzoGiuntiLavederLi2017}; sterile CP phases fixed to zero
  (no sensitivity in that data combination). The primary benchmark used
  throughout \cref{ch:sterile-results}.} \\[34pt]
\parbox[t]{3.4cm}{\code{sterile\_\allowbreak{}3p1\_\allowbreak{}benchmark\_\allowbreak{}icecube}} &
  \parbox[t]{9.2cm}{$\theta_{14}=0$, $\theta_{24}\approx9.22^\circ$
  ($\sin^22\theta_{24}=0.10$), $\theta_{34}=0$,
  $\Delta m^2_{41}=0.3\ \mathrm{eV^2}$. IceCube atmospheric $\nu_\mu$
  disappearance benchmark \parencite{IceCube2020SterileNumuDisappearance};
  $\theta_{14},\theta_{34}$ set to zero since IceCube atmospheric neutrinos
  are primarily sensitive to $\theta_{24}$.} \\[26pt]
\parbox[t]{3.4cm}{\code{sterile\_\allowbreak{}3p1\_\allowbreak{}two\_\allowbreak{}flavour\_\allowbreak{}limit}} &
  \parbox[t]{9.2cm}{$\theta_{12}=\theta_{13}=\theta_{23}=0$,
  $\theta_{14}=15^\circ$, $\theta_{24}=12^\circ$, $\theta_{34}=0$,
  $\Delta m^2_{41}=1.0\ \mathrm{eV^2}$. A mathematical construction, not a
  data-driven benchmark: switching off every SM active angle makes the
  $\nu_e$ disappearance channel reduce exactly to the textbook 2-flavour
  formula for any $\theta_{24}$. Used in \cref{sec:results-kinematics} to
  isolate the 2-flavour approximation from the active-sector suppression
  quantified there for the realistic presets above.} \\[40pt]
\bottomrule
\end{tabular}
\end{center}
\end{implbox}

\begin{refbox}
\cite{GariazzoGiuntiLavederLi2017,IceCube2020SterileNumuDisappearance,KoppMachadoMaltoniSchwetz2013}.
\end{refbox}

\section{Building an oscillation scenario with the sterile extension}

\begin{examplebox}[title={Attaching a sterile preset to a propagation config}]
\begin{lstlisting}[style=tpython]
from tpeanuts.config.propagation import PropagationConfig
from tpeanuts.util.context import RuntimeContext

context = RuntimeContext.resolve("cpu", torch.float64)
oscillation = PropagationConfig.oscillation_parameters_from_preset(
    "sterile_3p1_bestfit_giunti2017",   # any theta14_deg-carrying preset
    context=context,
)
# oscillation.pmns is now a PMNS_sterile (n_flavours == 4), and
# oscillation.mass_spectrum carries DeltamSq41; every downstream builder
# (hamiltonian_reduced/flavour, evolutor_numerical,
# evolutor_perturbative_segment, medium.*) generalises automatically.
\end{lstlisting}
\end{examplebox}

\begin{examplebox}[title={Enabling the sterile neutral-current term}]
\begin{lstlisting}[style=tpython]
from tpeanuts.medium.earth.probability import earth_probability_state

# include_matter_nc=None (default) auto-resolves to True here, since the
# preset is sterile and the profile below carries neutron-density data.
P = earth_probability_state(
    nustate, profile, oscillation, E_MeV, eta, depth_m,
    method="numerical", include_matter_nc=None,
)
\end{lstlisting}
\end{examplebox}

\section{Observables and simulations that can be built with the sterile extension}

\begin{implbox}
Every observable \texttt{tpeanuts} exposes for the Standard Model can be
re-evaluated with a \pyclass{PMNS\_sterile} in place of \pyclass{PMNS\_SM},
subject to the method-compatibility notes below:
\begin{itemize}
  \item \textbf{Vacuum $\nu_e$/$\nu_\mu$ disappearance.} The 4-flavour
    generalisation of the standard vacuum oscillation probability, used
    throughout \cref{ch:sterile-results} to compare against the textbook
    2-flavour formula and to locate short-baseline experiments in the
    $(L,E)$ plane.
  \item \textbf{Solar $P_{ee}(E)$.} \pyfunc{medium.solar.probability.solar\_probability\_mass}/
    \code{\_state} with \code{method="numerical"} or
    \code{method="adiabatic\_exact"} (any $N$, including $N=4$); the
    closed-form \code{method="adiabatic\_approximated"} two-level
    approximation remains 3-flavour-only, exactly as for NSI.
  \item \textbf{Earth regeneration.}
    \pyfunc{medium.earth.probability.earth\_probability\_state}/
    \code{\_numerical}/\code{\_analytical}, with the sterile neutral-current
    term (\S\ref{sec:sterile-matter}) available via \code{include\_matter\_nc}
    whenever the Earth profile was built with neutron-density coefficients.
  \item \textbf{Atmospheric propagation.} \pyfunc{medium.atmosphere.evolutor.atmosphere\_evolutor}
    and \pyfunc{atmosphere\_probability\_transition}, for $\nu_\mu$
    disappearance oscillograms directly comparable to the IceCube DeepCore
    benchmark.
  \item \textbf{Combined with NSI.} Any of the above with \code{oscillation.nsi}
    also set (companion \texttt{BSM\_NSI} document) -- the two BSM
    extensions compose freely, as documented in \cref{sec:impl-kinetic}'s
    case~3.
\end{itemize}
\end{implbox}

\begin{notebox}
As for NSI, \code{method="adiabatic\_approximated"} (the closed-form
two-level solar approximation and its Landau-Zener correction) rejects any
BSM extension outright, including the sterile one, since its
matter-mixing-angle formulas have no $N=4$ generalisation.
\pyfunc{analytic\_eigenvalues=True} (\cref{sec:impl-ferrari}) is available
for both the perturbative Earth/atmosphere paths at $N=4$, but is
incompatible with \code{method="numerical"}, which never diagonalises the
Hamiltonian at all.
\end{notebox}
